\section{Organización del trabajo en
grupo}\label{organizaciuxf3n-del-trabajo-en-grupo}

\subsection{Partes que aborda el TFE, distribución y estructura de la
memoria}\label{partes-que-aborda-el-tfe-distribuciuxf3n-y-estructura-de-la-memoria}

Organización del trabajo en grupo - Desarrollo de la memoria

{\def\LTcaptype{none} % do not increment counter
\begin{longtable}[]{@{}
  >{\raggedright\arraybackslash}p{(\linewidth - 2\tabcolsep) * \real{0.5000}}
  >{\raggedright\arraybackslash}p{(\linewidth - 2\tabcolsep) * \real{0.5000}}@{}}
\toprule\noalign{}
\begin{minipage}[b]{\linewidth}\raggedright
Apartado de la memoria
\end{minipage} & \begin{minipage}[b]{\linewidth}\raggedright
Responsables
\end{minipage} \\
\midrule\noalign{}
\endhead
\bottomrule\noalign{}
\endlastfoot
Introducción & Alumno 1, Alumno 2 y Alumno 3 \\
Estado del arte & Alumno 1, Alumno 2 y Alumno 3 \\
Objetivos y metodología de trabajo & Alumno 1, Alumno 2 y Alumno 3 \\
Desarrollo específico de la contribución: parte individual 1 & Alumno
1 \\
Desarrollo específico de la contribución: parte individual 2 & Alumno
2 \\
Desarrollo específico de la contribución: parte individual 3 & Alumno
3 \\
Desarrollo específico de la contribución: parte 4 & Alumno 1, Alumno 2 y
Alumno 3 \\
Conclusiones y trabajo futuro & Alumno 1, Alumno 2 y Alumno 3 \\
\end{longtable}
}

\subsection{Mecanismos de coordinación
empleados}\label{mecanismos-de-coordinaciuxf3n-empleados}

Con el objetivo de garantizar un seguimiento continuo del trabajo y
mantener la coordinación entre los miembros del grupo, se estableció una
reunión interna semanal en la que cada integrante exponía las tareas
realizadas durante el periodo correspondiente y las posibles
dificultades encontradas. Estas reuniones permitían, además, debatir las
siguientes líneas de trabajo y planificar las tareas a desarrollar en
las semanas posteriores.

Como herramientas de comunicación y colaboración se utilizaron
principalmente WhatsApp para la comunicación diaria, Microsoft Teams
para la realización de reuniones telemáticas, un repositorio compartido
en GitHub para la gestión y sincronización del código fuente, y un
documento compartido de Microsoft Word para la elaboración conjunta de
la memoria.

Adicionalmente, tras cada una de las entregas parciales previstas en la
planificación del Trabajo Fin de Máster, se mantuvo una reunión de
seguimiento con el director del trabajo. Estas sesiones permitieron
revisar el progreso realizado, recibir retroalimentación sobre los
resultados obtenidos y definir las acciones necesarias para las
siguientes fases del proyecto.
